%---
% slug: axosim-fly-geometry
% title: Finding Visual Dynamics in a Simulated Fly Brain
% category: research
% eyebrow: NEURAL SIMULATION
% summary: We run a pretrained AxoSim neuron model on a 45,669-cell fly connectome and find a generalizable motion-direction geometry in its T4/T5 voltage responses, alongside stable oscillatory behavior when exposing image stimuli.
% author: Axym Labs
% draft: false
% date: 2026-10-01
% rank: 1
%---
\documentclass[11pt]{article}
\usepackage{publication-web}
\begin{document}

\publicationtitle{Finding Visual Dynamics in a Simulated Fly Brain}
\publicationauthors{Axym Labs}
\publicationaffiliations{Darmstadt, Germany}

\begin{publicationhead}
\makepublicationtitle
\end{publicationhead}

\ifdefined\HCode
  \ifvmode\IgnorePar\fi\EndP
  \HCode{<figure class="publication-figure"><video class="publication-video" controls muted playsinline preload="metadata" poster="/work/axosim/axosim-neural-landscape-poster.png"><source src="/work/axosim/axosim-neural-landscape.mp4" type="video/mp4">Your browser does not support embedded video.</video><figcaption>Neural-state trajectory under naturalistic and Voyager images.</figcaption></figure>}\ShowPar
\else
  \publicationfigure[\linewidth]
    {figures/axosim-neural-landscape-poster.png}{/work/axosim/axosim-neural-landscape-poster.png}
    {A natural image beside a three-dimensional neural-state trajectory.}
    {Neural-state trajectory under naturalistic and Voyager images.}
\fi

In this research update, we examine how a simulated visual circuit of a fly represents natural images and optic flow. Various demonstrations of learningo or embodyment of the male fruit fly's brain circulate on the internet, but almost all of them are not using the its neural network in a meaningful way, other than for purposes of virality. We instantiate each neuron in the fruit fly's connectome (neural graph) as a pretrained AxoSim--Lite model, perform dimensionality reduction, record membrane voltage from T4 and T5 cells, and test whether motion direction is encoded geometrically in a low-dimensional space \cite{axosim,flyvis}. On held-out images, a readout of the circuit's voltage response reaches 56.25\% accuracy over twelve directions and a 65.97\% correlation between neural distance and true circular angle distance.

Our results support an emerging picture of neural simulation as a source of inspectable population geometry: We observe that raw states are heavily influenced by image-specific texture, while variables regarding motion can be extracted from a simpler anatomically structured subspace. Under sustained visual stimuli, the fly brain model first descends towards a point in activation space, then oscillates stably around it. Across six constant images, full-state perturbation-and-return tests establish evidence for input-dependent limit cycles, replicating observed stable oscillatory population dynamics with AxoSim--Lite models trained on a different cell type. This article follows the structure of Goodfire's \textit{Finding the Tree of Life in Evo 2}, including its choices of visualization \cite{goodfire}. We release the code used for these experiments at \href{https://github.com/DavideWiest/axosim-demo}{github.com/DavideWiest/axosim-demo}.

\section{Background: Why do interpretability on scientific simulators?}

Scientific machine-learning models are becoming useful across molecular biology, materials science, weather, and neuroscience. Neuron surrogates have a distinctive opportunity: they can make large circuit simulations faster while retaining variables such as membrane voltage that have a direct interpretation and can connect observations to prior research. These models also have qualities that make interpretability both difficult and valuable:

\begin{itemize}
\tightlist
\item A surrogate can reproduce single-neuron traces while still changing the collective dynamics of a large connected circuit. Population-level interpretation can reveal those changes.
\item A connectome-constrained model can carry biological structure and model-specific artifacts at the same time. Interpretability helps distinguish a signal supported by the graph from one introduced by preprocessing or analysis.
\item Neural inputs and internal states are high-dimensional and temporally coupled, necessitating modern dimensionality reduction techniques.
\end{itemize}

At Axym Labs, we are building neuron models and the tooling needed to test them inside large biological circuits. As part of that work, we study how the resulting population states organize under controlled stimuli---in this case, visual motion across natural images and sustained visual input in a released fly visual-system graph.

The FlyVis connectome contains 45,669 cells and 1,513,231 edges. We use AxoSim's inference platform and instantiate each neuron as one small AxoSim--Lite model. Every result here comes from its simulated voltage rather than a hand-authored activity trace, gait controller, steering rule, or behavioral replay (as recent viral demonstrations related to fly connectomes do) \cite{flybraindemo}. However, AxoSim--Lite was trained as pyramidal L5 neuron surrogates, which means the following analysis does not capture the effect of different cell types in the population.

\section{Understanding model ontologies: how does the simulated circuit represent visual motion?}

The direction of optic flow, i.e. the direction at which an image moves across the fly's visual field, and angular distance, the minimal angle between two such directions are particularly interesting objects to analyze geometrically.

Firstly, because have approximate ground truth for these variables, such that we can measure the networks's organization against ground truth angular distances between image frames. Each photograph moves at one of twelve angles separated by 30 degrees, and circular angular distance defines the relationship between every pair of labels. Complete photographs, rather than individual trials, are held out to prevent the decoder from recognizing image texture. Secondly, because T4/T5 cells are known to respond selectively to motion direction. 

This article addresses three questions: Can the circuit distinguish motion direction across natural images withheld from analysis selection? Do nearby directions occupy nearby regions of neural state space? What happens to the population trajectory when a visual input is held in place long enough for adaptation to develop?

The display set contains 500 ImageNet validation images and 500 licensed iNaturalist observations \cite{imagenet,inaturalist}. Their WordNet and taxonomic labels provide a useful, arguably canoncial way to organize the image cohort, which we can use to compare observations to. This metadata is not used to fit and models or estimators. 

\publicationfigure[0.92\linewidth]
  {figures/taxonomic-tree.png}{/work/axosim/taxonomic-tree.png}
  {A radial tree organizing the ImageNet and iNaturalist labels in the display cohort.}
  {Image-label hierarchy of the display cohort.}

\section{Dataset construction: disentangling motion from image texture}

All 12,000 trials fit the unsupervised UMAPs, and 720 trials from the generalization metrics are measured on a separate 60-image cohort. UMAPs of the full recording show substantial structure.

\publicationfigure[\linewidth]
  {figures/activation-umaps.jpg}{/work/axosim/activation-umaps.jpg}
  {Three UMAP views of the same AxoSim voltage states, colored by source, label group, and photograph.}
  {Raw voltage UMAPs by source, label group, and photograph.}

However, the direct embedding of visual-neuron voltage is dominated by variables such as edges, contrast, and spatial frequency. Since we want to test specifically for motion ditsance, we remove that confound from the scored analysis by translating every photograph in all twelve directions, averaging voltage over the final 120 milliseconds, and subtracting each photograph's angle-averaged response.

\publicationfigure[\linewidth]
  {figures/sampling-and-averaging.png}{/work/axosim/sampling-and-averaging.png}
  {A natural photograph moves across the retinal lattice before the final voltage window is averaged.}
  {Stimulus sampling and voltage averaging.}
  
Specifically, we construct the dataset as follows:
  
\begin{itemize}
\tightlist
\item Center-crop each photograph to square luminance and translate it by 48 pixels in each of twelve directions.
\item Sample the moving image on the 721-column hexagonal retinal lattice and run the full 45,669-cell graph for 320 milliseconds.
\item Record all 5,768 T4a--d and T5a--d cells, inverse-transform AxoSim output channel 1 to soma millivolts, and average the last 120 milliseconds.
\item Keep 48 complete photographs for representation development and open twelve additional photographs only for the final score. The other 940 images affect unsupervised display geometry only.
\end{itemize}

Visual networks trained specifically for motion estimation would likely show stronger direction structure in their raw states.

\section{Experiments}

\subsection{Finding a visual motion manifold}

Recent interpretability work has identified circular manifolds for calendar features and colors, and helical manifolds for numbers and years \cite{engels2025,modell2025,kantamneni2025}. These examples show how meaning can be encoded in distributed, geometric relationships between population states, exposed as pairwise distances, cosine similarities or membership to clusters, for example.

Motion direction, described by an angle, has the same circular topology, and we can fortunately use true distances taken from the data, which allows us to analyze whether shortest paths through the voltage manifold follow a known physical variable.

Curved manifolds are often high-dimensional, and dimensions hidden by low-dimensional visualizations might separate nearby conditions, carry image-specific information, or reflect nuisance variation. In the tested circuit, we might expect color or texture to organize the data strongly because every direction is observed through a different retinal pattern.

Distances along the curved manifold may still be meaningful \cite{modell2025}. We therefore compare raw cosine and graph-geodesic distance with circular motion-angle separation before learning any supervised projection.

\publicationfigure[\linewidth]
  {figures/graph-construction.png}{/work/axosim/graph-construction.png}
  {Angular distances form a nearest-neighbor graph whose shortest paths define manifold distance.}
  {Nearest-neighbor graph for manifold distances.}

We follow the same basic analysis pattern used for other feature manifolds:

\begin{itemize}
\tightlist
\item Compare cosine similarity between centered voltage vectors with circular angular distance.
\item Construct the smallest connected nearest-neighbor graph under angular distance.
\item Compute geodesic distance by summing edge lengths along the shortest path between two states.
\item Evaluate the relationship only on complete-image splits so that repeated texture cannot cross the boundary.
\end{itemize}

We can't infer from activity directly: Cosine distance and connected-graph geodesic distance are effectively unrelated to motion angle, with correlations of -0.013 and 0.0069. The natural-image manifold is real, but its dominant geometry does not recover the variable we care about.

\publicationfigure[\linewidth]
  {figures/manifold-distance.jpg}{/work/axosim/manifold-distance.jpg}
  {Raw cosine and graph distance against circular motion-angle separation.}
  {Raw neural distance versus motion-angle separation.}

The largest deviations occur because texture-specific retinal responses pull states apart even when their motion directions match. This explains both the image clusters in the UMAP and the flat relationship in the distance plot.

Looking towards neurology, the neurons of the classes T4 and T5 provide a plausible route to the hidden signal: Each class tiles retinal position and carries direction-selective structure, so spatial summaries within each class can discard exact texture location while retaining how strongly the population responds.

This is a good example in interpretability of how structure alone is not evidence that the model represents the intended concept in a readily measureable way, or at all.

\subsection{Finding a flat representation}

Consider an analogous analysis of place representations. A curved embedding may contain culture, language, and visual appearance, while a two-dimensional latitude-longitude subspace preserves the distance of interest directly. For motion direction in particular, using a circular coordinate system is preferrable.

The raw voltage manifold contains extra information about each photograph. We therefore ask whether a low-dimensional linear transformation can expose a flatter direction representation underneath that curvature.

We learn a ten-dimensional transformation in which pairwise neural distance predicts circular angular separation while a reconstruction term retains the original centered representation. In simplified form,

\[
z = W(x-b), \qquad
\widehat d_{12}=\beta\,\arccos\!\left(\frac{z_1^\top z_2}{\lVert z_1\rVert\lVert z_2\rVert}\right).
\]

The reconstruction term makes the variance retained by the subspace measurable and prevents the distance objective from collapsing onto an arbitrary low-rank solution. A development-only sweep compared anatomically fixed representations before the final photographs were opened.

Cross-validation holds out complete photographs. Six grouped folds select among spatial means and standard deviations, quantiles, moments, class means, and a PCA over all 5,768 retinal positions. The winning representation---mean and standard deviation within each T4/T5 class---is then frozen.

The fitted ten-dimensional space reaches a 0.7121 distance correlation on development-image pairs and 0.6597 on final-image pairs while retaining 0.8939 of the selected representation's variance. A classifier trained on the 48 development photographs predicts one of twelve directions on the twelve final photographs with 0.5625 balanced accuracy, compared with 0.0833 chance.

\publicationfigure[\linewidth]
  {figures/subspace-distance.jpg}{/work/axosim/subspace-distance.jpg}
  {Distance correlations in development, held-out, and cross-split image pairs.}
  {Motion-subspace distance correlations across image splits.}

The gap between development and final distance correlation is modest, but the experiment is still small: the final score contains twelve photographs and 144 trials. The image-bootstrap interval for accuracy is 0.4653 to 0.6597.

These results support a flat-structure-plus-deviations picture. A compact direction coordinate explains a large part of the anatomically summarized voltage, while the raw retinal state retains additional image-specific curvature.

\publicationfigure[0.62\linewidth]
  {figures/subspace-umap.png}{/work/axosim/subspace-umap.png}
  {A two-dimensional UMAP of the learned motion subspace colored by direction.}
  {Motion-subspace UMAP by direction.}

\publicationfigure[0.86\linewidth]
  {figures/learned-3d-umap.png}{/work/axosim/learned-3d-umap.png}
  {A fixed-camera three-dimensional view of the learned motion subspace.}
  {Three-dimensional UMAP of the motion subspace.}

\subsection{Images as visual ``styles''}

We have examined the direction structure, but what signals produce the remaining curvature?

The voltage vectors can carry at least two types of information:

\begin{itemize}
\tightlist
\item The general visual ``style'' of an image, including contrast, edge density, spatial frequency, and the distribution of luminance over the retinal lattice.
\item Specific local motifs that identify a photograph, such as a branch, horizon, façade, or high-contrast boundary at a particular retinal position.
\end{itemize}

Natural images contain many low-level statistics that can organize a representation meaningfully without proper object recognition, which most likely is a problem too hard for the fly circuit. Mean luminance, RMS contrast, and edge energy are especially relevant because the visual graph receives luminance sampled at fixed retinotopic coordinates.

We find preliminary evidence for this account. Photograph identity dominates raw-state neighborhoods, while the largest residual component and low-level image statistics align with visible directions through the untransformed geometry. Removing the angle-averaged image response and pooling within named cell classes exposes direction more clearly.

\publicationfigure[\linewidth]
  {figures/subspace-deviations.jpg}{/work/axosim/subspace-deviations.jpg}
  {Raw voltage geometry colored by its largest deviation from the learned motion subspace.}
  {Raw voltage geometry colored by subspace deviation.}

\publicationfigure[\linewidth]
  {figures/image-statistics.jpg}{/work/axosim/image-statistics.jpg}
  {Associations between the neural subspace and luminance, contrast, and edge energy.}
  {Neural geometry and low-level image statistics.}

\publicationfigure[\linewidth]
  {figures/retinotopic-vision-manifold.png}{/work/axosim/retinotopic-vision-manifold.png}
  {Natural images, measured retinotopic voltage fields, and the shared motion geometry.}
  {Natural images, retinal voltage responses, and motion geometry.}

Taken together, these figures separate two scales of the response. Retinotopic voltage preserves the spatial detail that distinguishes individual images, while pooling within T4/T5 classes exposes the lower-dimensional motion signal shared across them.

\section{Testing the oscillatory attractor}

The film at the top provides a temporal view of the same idea. Each pale point is one complete 2,884-neuron population state; the bright causal trail follows one continuous AxoSim response across twelve naturalistic images and twelve public-domain images carried on the Voyager Golden Record \cite{voyagerimages}. Across 23 cross-image transitions, the median distance to the target image's late recurrent orbit falls by a factor of 2.17. The dominant frequency stays the same between consecutive 1.5-second windows of the steady portion of all 24 intervals whereas amplitudes differ, ranging from $0.909\times$ to $1.302\times$. The population descends towards a point in activation space, then oscillates stably around it. This connects to observed limit-cycle attractors from work on recurrent neural population dynamics, such as those reviewed by Miller, alongside stimulus-linked 20--30 Hz central-brain oscillations and coherent oscillations of membrane potential recorded across neuron pairs in \textit{Drosophila} \cite{miller2016,grabowska2020,franco2025}.

To test whether this phenomenon is replicated precisely rather than just conceptually, we hold six images constant for four seconds and then displace every free recurrent-state variable with Gaussian noise scaled to either $0.10$ or $0.25$ times its state component's standard deviation around the orbit. The state comprises eight AxoSim hidden values and four propagated activity values for each of 39,901 recurrent cells. We exclude the 5,768 retinal-input cells because their propagated activity is externally clamped before every graph update, so their hidden states cannot affect the recurrent dynamics. At 20-millisecond intervals, we measure the minimum standardized RMS distance from the perturbed state to any phase of the pre-perturbation orbit and subtract the distance of a matched unperturbed continuation.

Both preregistered perturbation amplitudes pass the return gate. At $0.10\sigma$, the median late excess reaches the matched-control floor, with a 95\% scene-cluster bootstrap upper bound of 0.0133, and 29 of 30 trials return for at least 500 milliseconds by 2.5 seconds. At $0.25\sigma$, the median also reaches the matched-control floor, the bootstrap upper bound is 0.0395, and 27 of 30 trials meet the same sustained-return criterion. The four failures occur on one city trial at the lower amplitude and on two apple trials as well as one Voyager Earth trial at the higher amplitude.

\publicationfigure[\linewidth]
  {figures/limit-cycle-return.png}{/work/axosim/limit-cycle-return.png}
  {Full-state distance to the unperturbed orbit after two perturbation amplitudes, with late return ratios for six constant images.}
  {Full-state perturbations return toward the unperturbed orbit.}

The phase test is measured on the full 478,812-dimensional free recurrent state. Together with contraction at both perturbation amplitudes, this establishes evidence for input-dependent limit cycles across the six tested images.

\section{Discussion}

In this research update, we use controlled motion as an external ruler for a fly visual circuit. The experiment separates a texture-dominated raw manifold from a learned ten-dimensional direction representation, and the film shows how the same AxoSim population moves through activation space under sustained image stimuli.

Although preliminary, the results add evidence to two useful observations about neural population geometry:

\begin{itemize}
\tightlist
\item A simple coordinate transformation can expose stimulus distance tracking in the analyzed fly circuit, while retaining the bulk of the variance.
\item After motion direction is isolated in this compact subspace, the remaining variation reflects image-specific texture and low-level visual statistics.
\item Full-state perturbation-and-return tests establish evidence for input-dependent limit cycles across six images, replicating observed stable oscillatory population dynamics with AxoSim--Lite models trained on mammalian layer-5 pyramidal neurons.
\end{itemize}

Population manifolds appear useful, but methods that treat neural features as isolated directions do not capture their geometry well \cite{modell2025,goodfire}. Here we use a supervised distance objective to expose the motion subspace, alongside unsupervised displays of the raw voltage. Developing methods that recover such structure without a chosen external metric, and with less manual analysis, would make the approach more scalable. Manifolds and attractors should not be inferred from a clean visualization alone: complete-image holdouts, a label-preserving permutation test, and uncertainty over images support the motion claim, while a continuous mean-gray control tests the film's temporal separation. Here, the full-state perturbation test establishes evidence for attraction to input-dependent oscillatory orbits in the recurrent state.

These results suggest a practical path for interpreting scientific neural models: choose a variable with an external metric, isolate it from obvious input confounds, test generalization across the confounding unit, and only then inspect the remaining geometry. Better population-level understanding is a starting point for auditing the surrogate, comparing it with neural recordings, and eventually connecting validated sensory dynamics to a neuron-to-muscle pathway without inserting an engineered controller.

\begin{thebibliography}{99}
\bibitem{axosim} D. Wiest. \href{https://github.com/Axym-Labs/AxoSim-paper}{AxoSim: Learned Neuron Models for Fast Population Simulation and Inference-Time Adaptation}. Axym Labs technical report, 2026.
\bibitem{flyvis} J. Lappalainen et al. \href{https://doi.org/10.1038/s41586-024-07939-3}{Connectome-constrained networks predict neural activity across the fly visual system}. \textit{Nature}, 2024.
\bibitem{goodfire} M. Pearce et al. \href{https://www.goodfire.com/research/phylogeny-manifold}{Finding the Tree of Life in Evo 2}. Goodfire, 2025.
\bibitem{flybraindemo} J. Paquette. \href{https://github.com/jppaquet/flybrain}{flybrain: Simulate the Complete Connectome of the Male Fruit Fly}. GitHub repository, 2026.
\bibitem{imagenet} J. Deng, W. Dong, R. Socher, L.-J. Li, K. Li, and L. Fei-Fei. \href{https://doi.org/10.1109/CVPR.2009.5206848}{ImageNet: A Large-Scale Hierarchical Image Database}. \textit{CVPR}, 2009.
\bibitem{inaturalist} G. Van Horn et al. \href{https://doi.org/10.1109/CVPR.2018.00914}{The iNaturalist Species Classification and Detection Dataset}. \textit{CVPR}, 2018.
\bibitem{engels2025} J. Engels, E. J. Michaud, I. Liao, W. Gurnee, and M. Tegmark. \href{https://arxiv.org/abs/2405.14860}{Not All Language Model Features Are One-Dimensionally Linear}. \textit{ICLR}, 2025.
\bibitem{modell2025} A. Modell, P. Rubin-Delanchy, and N. Whiteley. \href{https://arxiv.org/abs/2505.18235}{The Origins of Representation Manifolds in Large Language Models}. arXiv:2505.18235, 2025.
\bibitem{kantamneni2025} S. Kantamneni and M. Tegmark. \href{https://arxiv.org/abs/2502.00873}{Language Models Use Trigonometry to Do Addition}. arXiv:2502.00873, 2025.
\bibitem{voyagerimages} NASA Science. \href{https://science.nasa.gov/mission/voyager/golden-record-contents/images/}{Images on the Golden Record}.
\bibitem{miller2016} P. Miller. \href{https://doi.org/10.12688/f1000research.7698.1}{Dynamical systems, attractors, and neural circuits}. \textit{F1000Research}, 2016.
\bibitem{grabowska2020} M. J. Grabowska, R. Jeans, J. Steeves, and B. van Swinderen. \href{https://doi.org/10.1073/pnas.2010749117}{Oscillations in the central brain of Drosophila are phase locked to attended visual features}. \textit{PNAS}, 2020.
\bibitem{franco2025} M. G. Franco et al. \href{https://doi.org/10.1016/j.isci.2025.112942}{Neuronal synchronization in Drosophila}. \textit{iScience}, 2025.
\bibitem{flyvis} J. Lappalainen et al. \href{https://doi.org/10.1038/s41586-024-07939-3}{Connectome-constrained networks predict neural activity across the fly visual system}. \textit{Nature}, 2024.
\bibitem{goodfire} M. Pearce et al. \href{https://www.goodfire.com/research/phylogeny-manifold}{Finding the Tree of Life in Evo 2}. Goodfire, 2025.
\bibitem{voyagerimages} NASA Science. \href{https://science.nasa.gov/mission/voyager/golden-record-contents/images/}{Images on the Golden Record}.
\end{thebibliography}

\end{document}
